% SPDX-License-Identifier: CC-BY-4.0
\documentclass[11pt,letterpaper]{article}

\usepackage[T1]{fontenc}
\usepackage{amsmath,amssymb}
\usepackage{tikz}
\usepackage{pgfplots}
\pgfplotsset{compat=1.18}
\usepackage{algorithm}
\usepackage{algpseudocode}
\usepackage{booktabs}
\usepackage{enumitem}
\usepackage[hidelinks]{hyperref}

\newcommand{\N}{\mathbb{N}}
\newcommand{\Z}{\mathbb{Z}}
\newcommand{\Q}{\mathbb{Q}}
\newcommand{\R}{\mathbb{R}}
\newcommand{\C}{\mathbb{C}}
\newcommand{\Map}[3]{#1 \colon #2 \to #3}
\newcommand{\Mod}[1]{#1\text{-}\mathrm{Mod}}
\newcommand{\GroupHom}{\mathrm{Hom}_{\mathbf{Grp}}}
\newcommand{\RingHom}{\mathrm{Hom}_{\mathbf{Ring}}}

\newtheorem{theorem}{Theorem}
\newtheorem{lemma}[theorem]{Lemma}
\newtheorem{corollary}[theorem]{Corollary}
\newtheorem{definition}[theorem]{Definition}
\newtheorem{remark}[theorem]{Remark}
\newenvironment{proof}{\noindent\textit{Proof.}}{\hfill\(\square\)}

\title{Generalized Collatz Library}
\author{}
\date{\today}

\begin{document}
\maketitle
\section*{Unit homogeneous raw maps}

\begin{abstract}
Let $D\geq2$.  The unit-multiplier homogeneous raw map divides a positive
integer by $D$ when possible and otherwise leaves it fixed.  Writing the seed
uniquely as $n=rD^k$ with $D\nmid r$ gives the exact raw-time formula
$T_D^j(n)=rD^{k-j}$, where subtraction is truncated at zero.  Thus the orbit
first reaches $r$ at time $k$, remains there, and never exceeds its seed.  The
argument classifies all positive fixed points and cycles and includes an
arbitrary-precision executable certificate.
\end{abstract}

\section{Raw map and canonical data}

\begin{definition}[Unit homogeneous map]
For an integer $D\geq2$, define
\[
 T_D(n)=
 \begin{cases}
  n/D,&D\mid n,\\
  n,&D\nmid n
 \end{cases}
 \qquad(n\in\N_{>0}).
\]
The divisibility test precedes the two branch formulas.  In particular, if
the map is extended to zero, then zero takes the division branch and remains
zero.
\end{definition}

\begin{lemma}[Canonical whole-base decomposition]
\label{lem:decomposition}
For every $n\in\N_{>0}$ there is a unique pair $(r,k)$ such that
\[
 n=rD^k,\qquad r\in\N_{>0},\qquad k\in\N,
 \qquad D\nmid r.
\]
\end{lemma}

\begin{proof}
The nonempty finite set
\[
 \{e\in\N:D^e\mid n\}
\]
has a largest member $k$: indeed $D^e\geq2^e>n$ for every sufficiently large
$e$.  Put $r=n/D^k$.  Then $n=rD^k$, and $D\mid r$ would imply
$D^{k+1}\mid n$, contrary to maximality.  Positivity of $r$ follows from
positivity of $n$.

For uniqueness, suppose
$rD^k=sD^\ell$ with $D\nmid r$ and $D\nmid s$.  If $k<\ell$, cancellation of
$D^k$ gives $r=sD^{\ell-k}$, so $D\mid r$, a contradiction.  If $\ell<k$,
the symmetric argument gives $D\mid s$.  Hence $k=\ell$, and cancellation
gives $r=s$.
\end{proof}

\begin{definition}[Entry certificate]
For $n=rD^k$ as in Lemma~\ref{lem:decomposition}, call $r$ the
\emph{canonical fixed factor} and $k$ the \emph{entry time}.
\end{definition}

\section{Exact raw-time classification}

\begin{lemma}[One-step exponent update]
\label{lem:one-step}
If $D\nmid r$, then for every $e\in\N$,
\[
 T_D(rD^e)=rD^{e-1},
\]
where $e-1$ is truncated at zero.
\end{lemma}

\begin{proof}
If $e=0$, then $rD^e=r$ is not divisible by $D$, so the fixed branch gives
$r$.  If $e>0$, then $D\mid rD^e$, and the division branch gives
$rD^e/D=rD^{e-1}$.
\end{proof}

\begin{theorem}[Exact iterate formula]
\label{thm:iterate}
Let $n=rD^k$ be the canonical decomposition of a positive seed.  For every
$j\in\N$,
\[
 T_D^j(n)=rD^{k-j},
\]
with truncated subtraction.
\end{theorem}

\begin{proof}
Induct on $j$.  At $j=0$ the formula is the decomposition of $n$.  If it
holds at $j$, Lemma~\ref{lem:one-step} gives
\[
 T_D^{j+1}(n)=T_D(rD^{k-j})=rD^{(k-j)-1}=rD^{k-(j+1)}.
\]
The final equality is the associativity of truncated subtraction on natural
numbers.
\end{proof}

\begin{theorem}[Exact entry and fixed tail]
\label{thm:entry}
For the same data, $T_D^k(n)=r$, and
\[
 T_D^{k+j}(n)=r
 \qquad\text{for every }j\in\N.
\]
If $0\leq j<k$, then $T_D^j(n)>r$.  Consequently $k$ is the first hitting
time of the canonical fixed factor.
\end{theorem}

\begin{proof}
Theorem~\ref{thm:iterate} at $j=k$ gives $rD^0=r$.  At time $k+j$, truncated
subtraction again gives exponent zero.  If $j<k$, then $k-j>0$.  Since
$D\geq2$ and $r>0$, one has $D^{k-j}>1$ and hence
$rD^{k-j}>r$.  This also proves that no earlier time reaches $r$.
\end{proof}

\begin{theorem}[Fixed points, cycles, and height]
\label{thm:classification}
For $n\in\N_{>0}$, the following are equivalent:
\[
 T_D(n)=n,
 \qquad D\nmid n.
\]
Every positive orbit is eventually fixed, every positive periodic orbit is
already fixed, and
\[
 T_D^j(n)\leq n
 \qquad\text{for all }j\in\N.
\]
Thus the positive cycles are exactly the one-point cycles supported on the
positive integers not divisible by $D$.
\end{theorem}

\begin{proof}
If $D\nmid n$, the second branch fixes $n$.  If $D\mid n$, then
$T_D(n)=n/D<n$ because $n>0$ and $D\geq2$, so $n$ is not fixed.
Theorem~\ref{thm:entry} proves eventual fixation.  A periodic orbit cannot
contain a strict decrease and later return, because each step is either a
division by $D$ or equality.  It is therefore fixed at every step.  Finally,
Theorem~\ref{thm:iterate} gives
$rD^{k-j}\leq rD^k=n$.
\end{proof}

\section{Reconstruction and computation}

For the positive homogeneous raw map with $q=1$, $c=0$, and $D\geq2$, the
canonical pair $(r,k)$ determines the entire orbit.
Lemma~\ref{lem:decomposition} proves that these data exist uniquely, and
Theorem~\ref{thm:iterate} reconstructs every raw iterate without changing the
clock.  Before entry, the remaining exponent $k-j$ is positive and drops by
one on each division step.  After exactly $k$ steps it is zero, and the
base-free state $r$ remains fixed.

The executable checker uses arbitrary-precision integers.  For each of the
$95$ divisors $2\leq D\leq96$ and all $4096$ seeds
$1\leq n\leq4096$, it constructs the canonical factor and entry time, replays
every strict pre-entry step, and verifies the fixed tail.  This is
\[
 95\cdot4096=389{,}120
\]
bounded certificates.  A separate seed $53\cdot37^{300}$ checks entry time
$300$ and fixed factor $53$.  These finite checks validate the implementation;
the universal conclusion is supplied by the proof above.

This result treats only $q=1$ and zero offset.  It does not classify positive
offsets, multipliers $q\geq2$, signed states, or divisors of absolute value
at least two with negative sign.  The positive map has infinitely many fixed
points, so the theorem does not assert a finite global attractor.

\end{document}
